\chapter{向量加法}

\begin{introduction}
  \item 一维 \texttt{mask}
  \item 二维 / 三维 stride
  \item 带宽绑定
  \item 文件：\texttt{01\_vector\_add/}
\end{introduction}

配套：\pyfile{kernel/triton/tutorial/01_vector_add/}。先跑：\pyfile{step01_add_1d.py} $\to$ \pyfile{step02_add_2d.py} $\to$ \pyfile{step03_add_3d.py}。

\section{公式与数据流}

\[
\mathrm{out}=x+y.
\]
每个 program 处理一块 \texttt{BLOCK} 元素；越界用 \texttt{mask}。二维、三维不要假设内存连续，用各自的 stride。二维 kernel 把 \texttt{pid} 拆成行块和列块：

\figtiletwo

\figbandwidth

\src{kernel/triton/tutorial/01_vector_add/step02_add_2d.py}{33}{47}{二维加法：stride 与二维 mask}

三维同理，再加一个 \texttt{pid} 轴。这是 launch 套路的操练场，不是算力测试。

\section{怎么跑}

\begin{runcmd}
\begin{lstlisting}[style=shell]
python kernel/triton/tutorial/01_vector_add/step01_add_1d.py
python kernel/triton/tutorial/01_vector_add/step02_add_2d.py
python kernel/triton/tutorial/01_vector_add/step03_add_3d.py
\end{lstlisting}
\end{runcmd}

\section{正确性}

\begin{note}
对照 \texttt{x + y}。elementwise 应对齐到 \texttt{atol} $10^{-5}$ 量级。
\end{note}

\section{效果}

\begin{remark}
这是\textbf{带宽绑定}：算的是一次加法，时间看 load/store。autotune 主要扫 \texttt{BLOCK} / \texttt{num\_warps}。不要拿它和 cuBLAS 比。
\end{remark}
